\chapter{Introdução}
\label{cap:introducao}

A avaliação quantitativa de políticas públicas enfrenta uma tensão permanente
entre duas perguntas aparentemente semelhantes, mas fundamentalmente distintas.
A primeira é descritiva: \textit{programas sociais estão associados a melhores
resultados?} A segunda é causal: \textit{esses programas causam melhores
resultados?} A diferença entre as duas não é apenas filosófica. Ela determina
se uma política deve ser expandida, reformulada ou descontinuada --- e,
portanto, como recursos públicos escassos são alocados.

Este guia é sobre a segunda pergunta. Seu objetivo é oferecer ao leitor as
ferramentas conceituais e práticas para estimar efeitos causais a partir de
dados observacionais e experimentais. Para isso, seguimos a estrutura que a
econometria moderna consagrou como o caminho rigoroso entre uma pergunta de
política e um número calculado em dados: \textbf{identificação},
\textbf{estimação} e \textbf{inferência}.


\section{O Problema e a Solução}

Toda análise causal começa com o \textbf{problema fundamental da inferência
causal}: para cada unidade de observação --- seja ela uma pessoa, um município
ou uma escola ---, só é possível observar o resultado sob \textit{um} dos
tratamentos possíveis. O salário de quem foi para a faculdade não pode ser
observado simultaneamente no mundo contrafactual em que a mesma pessoa não foi.
A nota de um aluno em turma pequena não pode ser comparada, para o mesmo aluno,
à nota que teria tirado em turma regular.

Essa limitação é estrutural, não uma falha dos dados. Ela implica que efeitos
causais individuais são intrinsecamente não-observáveis, e que a inferência
causal é, por natureza, um exercício de raciocínio sobre o que
\textit{não} foi observado.

A solução desenvolvida pela econometria moderna não elimina esse problema ---
ela o contorna por meio de hipóteses sobre o processo que gerou os dados. Se
indivíduos similares foram expostos a tratamentos distintos por razões que não
se relacionam com seus resultados potenciais, a comparação entre eles é
informativa sobre o efeito causal. Experimentos aleatórios garantem essa
condição por design. Métodos quasi-experimentais --- regressão com
descontinuidade, diferenças em diferenças, pareamento, variáveis instrumentais
--- procuram encontrá-la nos dados observacionais.


\section{A Revolução da Credibilidade}

A formação da econometria aplicada moderna foi marcada por uma virada
metodológica que ficou conhecida como \textbf{revolução da credibilidade}
\citep{Angrist2010}. Até os anos 1980, a estimação de efeitos causais dependia
amplamente de modelos estruturais com fortes restrições derivadas da teoria
econômica. Esses modelos produziam estimativas, mas sua credibilidade dependia
de hipóteses que raramente podiam ser verificadas nos dados.

A partir do trabalho seminal de pesquisadores como David Card, Joshua Angrist e
Guido Imbens, a prática econométrica mudou de direção: em vez de depender de
restrições teóricas implausíveis, passou-se a buscar \textit{variação exógena}
nos dados --- situações em que o tratamento é atribuído de forma suficientemente
aleatória para permitir identificação causal. A formalização do arcabouço de
resultados potenciais \citep{Rubin1974, Rubin1978, Neyman1923} forneceu o
vocabulário teórico unificado para essa abordagem.

Em 2021, o Prêmio Nobel de Ciências Econômicas reconheceu essa contribuição.
Card foi premiado pela aplicação de métodos quasi-experimentais a questões do
mercado de trabalho; Angrist e Imbens, pelas contribuições metodológicas à
análise de relações causais com dados observacionais. A Série Avaliação na
Prática, da qual este guia faz parte, é uma expressão direta desse legado.


\section{A Cadeia de Análise Causal}

O caminho que este guia percorre pode ser descrito como uma cadeia com
três elos:

\begin{enumerate}[leftmargin=*, label=\textbf{\arabic*.}]
  \item \textbf{Identificação} — Definir o parâmetro de interesse, derivar as
    hipóteses sobre o processo gerador dos dados que o tornam identificável, e
    especificar o estimando: a quantidade dos dados que, sob essas hipóteses,
    coincide com o parâmetro.

  \item \textbf{Estimação} — Construir um estimador que calcule o estimando a
    partir de uma amostra finita. Discutir suas propriedades (consistência,
    viés, eficiência) e as hipóteses funcionais adicionais que ele impõe.

  \item \textbf{Inferência} — Quantificar a incerteza amostral em torno da
    estimativa. Escolher estimadores de variância adequados à estrutura dos
    dados, realizar testes de hipótese, construir intervalos de confiança e
    verificar a plausibilidade das hipóteses de identificação por meio de
    testes de falsificação.
\end{enumerate}

Os três elos respondem a perguntas distintas e se articulam sequencialmente:
a inferência pressupõe a estimação, que pressupõe a identificação. Um erro
num dos elos não é corrigido pelos demais.

\begin{atencaobox}[Erro-padrão pequeno não é garantia de causalidade]
Um resultado com erro-padrão muito pequeno e p-valor minúsculo pode ser
completamente destituído de interpretação causal se a identificação for
inválida. A precisão estatística mede quão bem estimamos o estimando ---
não se o estimando corresponde ao parâmetro econômico de interesse.
Esta distinção percorre todo o guia.
\end{atencaobox}


\section{Exemplos Condutores}

Ao longo do guia, utilizamos dois exemplos para ilustrar os conceitos.

O primeiro é a questão dos \textbf{retornos à educação}: qual é o efeito
causal de um ano adicional de escolaridade sobre o salário? O viés de seleção
é imediato e intuitivo --- pessoas com mais educação tendem a ter maior
habilidade inata e condições socioeconômicas mais favorecidas, o que elevaria
seus salários mesmo sem o diploma. O exemplo tem longa tradição na literatura
\citep{Mincer1974, Card1994, Card1995} e dados abundantes no Brasil por meio
da PNAD Contínua.

O segundo exemplo, apresentado em detalhe no Capítulo~\ref{sec:star}, é o
\textbf{Experimento STAR} --- um experimento aleatório conduzido no Tennessee
entre 1985 e 1989, que avaliou o efeito do tamanho da turma sobre o desempenho
escolar \citep{Finn1990, Krueger1999}. O STAR serve como laboratório empírico
ideal para ilustrar identificação por aleatorização, estimação com efeitos
fixos, erros-padrão clusterizados, bootstrap, múltiplos testes, inferência
por permutação e testes de falsificação.


\section{Estrutura do Guia}

O Capítulo~\ref{sec:identificacao} apresenta o problema de identificação
causal: o arcabouço de resultados potenciais, os parâmetros de interesse
(ATE, ATT, ATU), o viés de seleção e as hipóteses de identificação.

O Capítulo~\ref{sec:estimacao} trata da passagem do estimando ao estimador:
o princípio plug-in, as propriedades dos estimadores, o papel da forma
funcional, o IPW, os estimadores de painel e o estimador de Lin.

O Capítulo~\ref{sec:inferencia} cobre inferência estatística: erros-padrão,
testes de hipótese, intervalos de confiança, clustering, bootstrap e
inferência por permutação, encerrando com os principais testes de
falsificação.

O Capítulo~\ref{sec:star} aplica todos esses elementos ao Experimento STAR.

O Capítulo~\ref{cap:conclusao} sintetiza as principais lições e conecta
este guia aos demais da Série Avaliação na Prática.

\begin{notabox}[Pré-requisitos]
Este guia pressupõe familiaridade com estatística básica (esperança,
variância, distribuição normal, testes de hipótese) e com o modelo de
regressão linear simples. Conhecimentos de cálculo são úteis mas não
indispensáveis. O código em R requer conhecimento básico da linguagem;
os pacotes utilizados são apresentados com exemplos autocontidos.
\end{notabox}